\section{공개 API 사용 가이드 (Public API)}

\subsection{핵심 헤더}
\begin{itemize}[leftmargin=*]
\item \texttt{<oqs/oqs.h>}: 통합 진입 헤더
\item \texttt{<oqs/common.h>}: 상태 코드, 초기화, 메모리, CPU 확장 체크
\item \texttt{<oqs/kem.h>}: KEM 객체/함수
\item \texttt{<oqs/rand.h>}: RNG 선택 및 바이트 생성
\end{itemize}

\subsection{라이프사이클 API}
\begin{longtable}{p{0.24\textwidth}p{0.70\textwidth}}
\toprule
함수 & 설명 \\
\midrule
\texttt{OQS\_init()} & 라이브러리 전역 초기화(OpenSSL/CPU feature 준비 포함) \\
\texttt{OQS\_destroy()} & 전역 자원 정리 \\
\texttt{OQS\_thread\_stop()} & 멀티스레드 환경에서 스레드 종료 전 정리 훅 \\
\texttt{OQS\_version()} & 라이브러리 버전 문자열 반환 \\
\bottomrule
\end{longtable}

\subsection{KEM API 흐름}
\begin{enumerate}[leftmargin=*]
\item \texttt{OQS\_KEM\_alg\_identifier(i)} 또는 문자열 상수로 알고리즘 선택
\item \texttt{OQS\_KEM\_new(name)}로 객체 생성
\item 길이 필드(\texttt{length\_public\_key} 등) 기반 버퍼 할당
\item \texttt{OQS\_KEM\_keypair}, \texttt{OQS\_KEM\_encaps}, \texttt{OQS\_KEM\_decaps} 호출
\item \texttt{OQS\_KEM\_free()} 및 민감정보 secure free
\end{enumerate}

\subsection{메모리 API 권장 패턴}
\begin{itemize}[leftmargin=*]
\item 민감정보(비밀키/공유비밀): \texttt{OQS\_MEM\_secure\_free}
\item 비민감정보(공개키/암호문): \texttt{OQS\_MEM\_insecure\_free}
\item 비교는 \texttt{memcmp} 대신 \texttt{OQS\_MEM\_secure\_bcmp} 우선 검토
\end{itemize}

\subsection{알고리즘 선택 시 주의}
\texttt{OQS\_KEM\_alg\_count()}와 \texttt{OQS\_KEM\_alg\_identifier()}로 열거되는 목록이
실행 시점의 enabled 상태와 동일하지 않을 수 있으므로,
실사용 시 \texttt{OQS\_KEM\_alg\_is\_enabled()} 체크를 병행한다.

\subsection*{Assumption}
\texttt{include/oqs/kem.h}의 \texttt{OQS\_KEM\_algs\_length}와 \texttt{src/kem/kem.c}의 식별자 배열 길이가
동일 전제를 가져야 안전하다. 코드 스냅샷상 두 값의 불일치 가능성이 보여 운영 전 점검이 필요하다.
